% Propietario conservado: /Users/ruben/Documents/New project/output/REVISION_ARGUMENTAL_APERTURAS_CIERRES_20260915/VII/delivery/MANUSCRIPT_VII_ES_EN_COMPLETE/source_es/manuscrito/80_observador_y_respuesta_relacional.tex
% Recuperación y articulación de los propietarios documentados en
% PROCEDENCIA_OBSERVADOR_RESPUESTA.md. Originales conservados íntegros.
% Antecedentes de integración: núcleo APP--TRIT--TPK, realización del
% transporte de borde, lectores areal/volumétrico y reloj energético.
% Las condiciones de representación se enuncian antes de cada uso.

\section{Observer, memory, and relational response}
\label{vii:vii:sec:observador-respuesta}

The reading operation acts on the transport published by the
APP--TRIT--TPK enriched state. Its domain preserves route, sheet, orientation,
carry and memory; the result accessible to the observer is a map
of that domain. This distinction allows three operations to be described
together: choosing an oriented representative of transport,
inscribing its record in an apparatus and the local response induced
by the global boundary. Each operation transmits different information.
The following construction specifies their maps and the conditions
allowing them to be composed.

\subsection{Boundary transport and intrinsic reading}
\label{vii:vii:sec:observador-ciclo}

The finite publication of transport provides an oriented cycle
\(C_n\), with \(n\geq1\), vertices \(0,\ldots,n-1\) and edges
\(e_j:j\longrightarrow j+1\pmod n\). We fix the lattice realization
\(L\) of the record, its vector space
\(W=L\otimes_{\mathbb Z}\mathbb R\), a positive inner product and
the nonzero primitive representative \(v\in L\) of the distinguished residual
class, in the declared chamber. The APP sum and product sheets
provide the transport polarizations; TRIT orientation and
TPK updating precede this linear publication. The cycle,
the lattice realization and the residual class are the data this
section receives from the transport constructed previously.

A \(0\)-cochain is a function on the vertices with values in
\(W\); a \(1\)-cochain assigns a vector to each oriented edge.
The coboundary acts by
\[
 (d\Phi)(e_j)=\Phi(j+1)-\Phi(j).
\]
We denote by \(\mathbf1_E\) the constant scalar cochain on the
edges. For a transport \(\omega\in C^1(C_n;W)\), its accumulated
publication \(\Gamma_\omega\) is obtained through partial sums and
affine interpolation; its closure defect is
\begin{equation}
 B(\omega)=\Gamma_\omega(1)-\Gamma_\omega(0)
 =\sum_{j=0}^{n-1}\omega(e_j).
 \label{vii:vii:eq:observador-defecto}
\end{equation}

\begin{theorem}[Exact decomposition and harmonic mode of the cycle]
\label{vii:vii:thm:observador-descomposicion}
Every transport has a unique decomposition
\[
 \omega=d\Phi+u\mathbf1_E,\qquad \Phi(0)=0,
 \qquad u=\frac1n\sum_j\omega(e_j).
\]
In particular, \(B(d\Phi)=0\) and \(B(\omega)=nu\).
\end{theorem}

\begin{proof}
Let \(y_j=\omega(e_j)-u\). The sum of the \(y_j\) is zero.
We define \(\Phi(k)=\sum_{j=0}^{k-1}y_j\); the value for \(k=n\)
coincides with \(\Phi(0)=0\), so the function is well defined
on the cycle, including its final edge. We have \(d\Phi(e_j)=y_j\).
If \(d\Phi+u\mathbf1_E=d\Psi+w\mathbf1_E\), the sum over edges gives
\(nu=nw\). Then \(d(\Phi-\Psi)=0\); connectivity of the cycle
and normalization at the initial vertex imply \(\Phi=\Psi\).
The identity \(\sum_j(\Phi(j+1)-\Phi(j))=0\) directly proves
telescopic cancellation.
\end{proof}

\begin{definition}[Admissible transport and normalized reading]
\label{vii:vii:def:observador-admisible}
Transport is admissible when its harmonic mode belongs to
\(\mathbb Rv\). On the space \(\mathcal T_{\rm adm}\) of those
transports we define the global reading and its density per phase:
\begin{equation}
 \mathfrak a_v(\omega)
 =\frac{\langle B(\omega),v\rangle}{\langle v,v\rangle},
 \qquad
 \mu_v(\omega)=\frac{\mathfrak a_v(\omega)}n.
 \label{vii:vii:eq:observador-lectura}
\end{equation}
We have \(B(\omega)=\mathfrak a_v(\omega)v\).
\end{definition}

\begin{theorem}[Determination of the intrinsic reading]
\label{vii:vii:thm:observador-unicidad}
The quotient
\(\mathcal T_{\rm adm}/dC^0(C_n;W)\) is one-dimensional, generated
by the class of \(v\mathbf1_E\). The map \(\mu_v\) is the
unique linear functional on \(\mathcal T_{\rm adm}\) satisfying
\[
 F(\omega+d\Phi)=F(\omega),\qquad F(v\mathbf1_E)=1.
\]
\end{theorem}

\begin{proof}
The preceding decomposition writes each admissible transport as
\(d\Phi+\lambda v\mathbf1_E\). Its class modulo exact cochains
is determined by \(\lambda\). The class of \(v\mathbf1_E\)
is nonzero because its sum is \(nv\neq0\). By linearity and
invariance, every functional with the indicated characteristics equals
\(F(\omega)=\lambda F(v\mathbf1_E)=\lambda\). The formula
\eqref{vii:vii:eq:observador-lectura} gives precisely \(\mu_v=\lambda\).
The global reading is \(\mathfrak a_v=n\mu_v\).
\end{proof}

\begin{definition}[Edge Observer]
\label{vii:vii:def:observador-borde}
Given the polarization of the published transport, a boundary observer
is a pair \(\mathcal O=(\varepsilon,\Psi)\), where
\(\varepsilon\in\{+1,-1\}\) and \(\Psi\in C^0(C_n;W)\).
Its reading action is
\begin{equation}
 \omega^{\mathcal O}=\varepsilon\omega+d\Psi.
 \label{vii:vii:eq:observador-accion}
\end{equation}
The sign fixes the orientation and the cochain fixes the trivialization.
Two polarization choices are compared through their published
transport; the relation \eqref{vii:vii:eq:observador-accion} describes
exactly the representatives differing by orientation and coboundary.
\end{definition}

\begin{proposition}[Covariance of the reading]
\label{vii:vii:prop:observador-covarianza}
The action preserves admissibility and satisfies
\[
 B(\omega^{\mathcal O})=\varepsilon B(\omega),\quad
 \mathfrak a_v(\omega^{\mathcal O})=\varepsilon\mathfrak a_v(\omega),
 \quad \mu_v(\omega^{\mathcal O})=\varepsilon\mu_v(\omega).
\]
Trivializations with the same orientation publish the same closure.
\end{proposition}

\begin{proof}
The sum of the coboundary \(d\Psi\) is zero. Therefore, the harmonic mode
transforms into \(\varepsilon u\), which still belongs to
\(\mathbb Rv\), and the defect transforms into \(\varepsilon B\).
The two scalar equalities are obtained by normalized projection.
\end{proof}

The unit of this reading is the residual representative \(v\), with its
internal normalization. Trivialization freedom acts on the
cochain representative; closure information resides in its
class. The description allows observers to be compared while keeping
genealogical transport, its linear publication and the scalar
evaluating the residual class separate.

\subsection{Boundary identity and cellular interior}
\label{vii:vii:sec:observador-stokes}

Let \(K\) be a finite oriented two-dimensional cell complex,
equipped with a fundamental \(2\)-chain \(c_K\) such that
\(\partial c_K=[C_n]\). This condition expresses that interior
incidences cancel with their multiplicities and orientations.
Let \(\Omega\in C^1(K;W)\) be an extension of \(\omega\) and
\(F_\Omega=d\Omega\) its additive discrete curvature.

\begin{theorem}[Cellular holographic identity of the reading]
\label{vii:vii:thm:observador-stokes}
With the preceding data,
\begin{equation}
 B(\omega)=\langle\Omega,\partial c_K\rangle
          =\langle d\Omega,c_K\rangle,
 \qquad
 \mathfrak a_v(\omega)
 =\frac{\langle\langle d\Omega,c_K\rangle,v\rangle}
        {\langle v,v\rangle}.
 \label{vii:vii:eq:observador-stokes}
\end{equation}
In particular, \(F_\Omega=0\) implies \(\mathfrak a_v(\omega)=0\).
If closure is nonzero, every extension of transport on that complex
has nonzero discrete curvature.
\end{theorem}

\begin{proof}
The pairing between chains and cochains satisfies
\(\langle d\Omega,c_K\rangle=\langle\Omega,\partial c_K\rangle\).
When it is expanded, each interior edge appears with the sum of the incidence
numbers of the adjoining cells; the condition on
\(\partial c_K\) leaves exactly the oriented boundary cycle.
The restriction of \(\Omega\) to that cycle is \(\omega\), whose sum
is \(B(\omega)\). Projection onto \(v\) gives the second formula
and the flatness implication.
\end{proof}

The equality compares two evaluations of the same extended cochain:
the accumulated boundary defect and the coboundary integrated over the interior.
Its curvature takes values in \(W\) and uses the additive cellular coboundary.
Realization by a geometric connection additionally preserves
its own coefficient space and transport law.

\subsection{Isometric inscription of the record in the apparatus}
\label{vii:vii:sec:observador-aparato}

Let \(\mathfrak M_n\) be a nonempty finite set of records
published at level \(n\), with result, route, carry, sheet and memory.
Its record space is
\(\mathcal H_n=\ell^2(\mathfrak M_n)\), with basis \(|m\rangle\).
The visible update \(F_n\) has a record lift
\[
 \widehat F_n(m)=(F_n(m),m).
\]
We take as records of the next level a set containing those
pairs. The second coordinate preserves the antecedent and makes the lift
injective. This construction specifies the memory retained by
inscription, including when the visible update merges values.

In a finite apparatus space \(\mathcal K_{A,n}\), we fix an
orthonormal family \((|a_m\rangle)_{m\in\mathfrak M_n}\) and define
\[
 \iota_n:\mathcal H_n\longrightarrow\mathcal K_{A,n},
 \quad \iota_n|m\rangle=|a_m\rangle,
 \qquad
 V_n:\mathcal H_n\longrightarrow\mathcal H_{n+1},
 \quad V_n|m\rangle=|\widehat F_n(m)\rangle.
\]
Let \(\mathcal P_n=\iota_n\mathcal H_n\) be the pointer subspace.

\begin{theorem}[Compatibility between updating and inscription]
\label{vii:vii:thm:observador-entrelazador}
The maps \(\iota_n\) and \(V_n\) are isometries. There exists a unique
isometry \(W_n^0:\mathcal P_n\longrightarrow\mathcal K_{A,n+1}\)
satisfying
\begin{equation}
 W_n^0\iota_n=\iota_{n+1}V_n.
 \label{vii:vii:eq:observador-cuadrado}
\end{equation}
If \(\dim\mathcal K_{A,n+1}\geq\dim\mathcal K_{A,n}\),
it extends to an isometry of all of \(\mathcal K_{A,n}\).
Through auxiliary complements equalizing dimensions, a unitary
extension between the enlarged spaces is obtained.
\end{theorem}

\begin{proof}
The basis images under \(\iota_n\) are orthonormal; the
images under \(V_n\) are also orthonormal by injectivity of
\(\widehat F_n\). For \(z=\iota_n\psi\) we define
\(W_n^0z=\iota_{n+1}V_n\psi\). The representation of \(z\) is unique,
and the three maps preserve its norm. The formula determines the map on
a basis of \(\mathcal P_n\) and proves the square.
If the target dimension is sufficient, the orthogonal complement of
\(W_n^0\mathcal P_n\) contains an orthonormal family of the size of
\(\mathcal P_n^\perp\); sending a basis of the latter to that family
constructs the extension. Finally both spaces are enlarged to
the same dimension and the two prescribed families are completed to orthonormal
bases. The map between those bases is unitary.
\end{proof}

When the record includes labels of a finite instrument, this same
inscription realizes its environment: if the operators
\(K_{y,a}:\mathcal H_S\to\mathcal H'_S\) satisfy
\(\sum_{y,a}K_{y,a}^*K_{y,a}=I\), the map
\[
 V_{\rm inst}\psi=\sum_{y,a}K_{y,a}\psi\otimes|y,a\rangle
\]
is isometric, since its squared norm is
\(\sum_{y,a}\langle\psi,K_{y,a}^*K_{y,a}\psi\rangle=\|\psi\|^2\).
The inscription \(\iota\) transports each label to its corresponding
pointer. The model thus fixes the premeasurement record; a
material device additionally requires its interaction and calibration.

\subsection{Accessible algebra and pointer complement}
\label{vii:vii:sec:observador-expectativa}

Define \(P_m=|a_m\rangle\langle a_m|\),
\(P=\sum_mP_m\), and \(P_\perp=I-P\). The algebra of the pointer
subspace is \(P\mathcal B(\mathcal K_{A,n})P\), whose unit is \(P\).
On it, the diagonal reader is
\[
 \mathcal D_P(A)=\sum_mP_mAP_m.
\]

\begin{proposition}[Record expectation and unital extension]
\label{vii:vii:prop:observador-pinzamiento}
The map \(\mathcal D_P\) is a completely positive,
trace-preserving and idempotent conditional expectation onto the
pointer algebra, with \(\mathcal D_P(P)=P\). Its image is
\(\bigoplus_m\mathbb CP_m\).
On the algebra of the complete apparatus, the map
\begin{equation}
 \mathcal E_P(A)=\sum_mP_mAP_m+P_\perp AP_\perp
 \label{vii:vii:eq:observador-complemento}
\end{equation}
is a unital, completely positive, trace-preserving
conditional expectation. Its image is
\[
 \left(\bigoplus_m\mathbb CP_m\right)
 \oplus\mathcal B(\mathcal P_n^\perp).
\]
Both maps are orthogonal projections for the Hilbert--Schmidt
inner product. Their spectra are contained in \(\{0,1\}\).
When eliminated components exist, the separation between the
fixed subspace and the kernel is one.
\end{proposition}

\begin{proof}
The projectors \(P_m\) are mutually orthogonal. In the complete
apparatus, the family that adds \(P_\perp\) sums to \(I\). Each term
\(A\mapsto P_mAP_m\) is completely positive: after any
matrix amplification it acts by congruence with \(I\otimes P_m\).
The same holds for \(P_\perp\). The sum preserves that property.
Cyclicity of the trace and the sum of the projectors prove
trace preservation on each domain. Orthogonality annihilates the
cross terms when the map is iterated, so it is idempotent.
Its fixed elements are precisely the stated blocks and the
map is bimodular with respect to that subalgebra.
Moreover,
\(\operatorname{Tr}(A^*P_mBP_m)
 =\operatorname{Tr}((P_mAP_m)^*B)\),
from which self-adjointness for Hilbert--Schmidt follows. A self-adjoint
projection has only the spectral values zero and one when
both subspaces are nonzero. If only the fixed subspace remains,
the map is the identity.
\end{proof}

The sum \(\sum_mP_mAP_m\), applied without the complement to the entire
apparatus, sends \(I\) to \(P\). Its natural unit belongs to the pointer
algebra. The formula \eqref{vii:vii:eq:observador-complemento} also preserves
the auxiliary block; an orthonormal resolution of that block
provides, when desired, diagonalization of the entire apparatus.

If \(U\) is a unitary transformation of the apparatus and
\(P'_m=UP_mU^*\), then
\[
 \mathcal E_{P'}(UAU^*)=U\mathcal E_P(A)U^*.
\]
The equality is obtained by substituting each projector into the sum.
Likewise, for \(A\in\mathcal B(\mathcal P_n)\), the square
\eqref{vii:vii:eq:observador-cuadrado} implies
\[
 \mathcal D_{P_{n+1}}(W_n^0A(W_n^0)^*)
 =W_n^0\mathcal D_{P_n}(A)(W_n^0)^*.
\]
Indeed, each matrix unit \(|a_m\rangle\langle a_{m'}|\)
is transported to the unit corresponding to the distinct records
\(\widehat F_n(m),\widehat F_n(m')\); the diagonal reader preserves
exactly the terms with \(m=m'\).

\subsection{Factorization of the response through the effective boundary}
\label{vii:vii:sec:observador-factorizacion}

Let \(\mathcal X_\infty^{\rm enr}\) be the domain of compatible enriched
states and let
\(b:\mathcal X_\infty^{\rm enr}\longrightarrow B_\infty\) be its global
boundary trace. We denote by \(B_{\rm ef}=\operatorname{im}b\)
the effectively published boundary. For a response space
\(\mathcal V_v\), we consider
\(I_v:\mathcal X_\infty^{\rm enr}\longrightarrow\mathcal V_v\).

\begin{theorem}[Relational response criterion]
\label{vii:vii:thm:observador-factorizacion}
There exists a unique map \(\overline I_v:B_{\rm ef}\to\mathcal V_v\)
such that \(I_v=\overline I_v\circ b\) if and only if
\begin{equation}
 b(x)=b(x')\quad\Longrightarrow\quad I_v(x)=I_v(x').
 \label{vii:vii:eq:observador-fibra}
\end{equation}
Uniqueness extends to all of \(B_\infty\) when \(b\) is
surjective.
\end{theorem}

\begin{proof}
A composition through \(b\) is constant on each fiber. Conversely,
for \(\beta\in B_{\rm ef}\), we define
\(\overline I_v(\beta)=I_v(x)\) with \(b(x)=\beta\).
Condition \eqref{vii:vii:eq:observador-fibra} guarantees independence
of the representative; every value of \(\overline I_v\) is determined
by some state of that fiber. If \(b\) is surjective,
\(B_{\rm ef}=B_\infty\). Otherwise, the composition equation
determines only the values on \(B_{\rm ef}\).
\end{proof}

This criterion is set-theoretic. For a continuous factorization, the
topologies and the quotient-map condition
of \(b:\mathcal X_\infty^{\rm enr}\to B_{\rm ef}\) are also preserved; with it,
continuity of \(I_v\) is equivalent to that of its factor. This specification
separates determination of a reading from its regularity properties.

Let \(\ell_v:\mathcal X_\infty^{\rm enr}\to\mathcal L_v\) be the
restriction accessible to a local neighborhood. The effective global character
is verified through two states satisfying
\begin{equation}
 \ell_v(x)=\ell_v(x'),\qquad I_v(x)\neq I_v(x').
 \label{vii:vii:eq:observador-testigo-global}
\end{equation}
If \(I_v\) factors through \(b\), these two states have distinct
boundaries. Moreover, the response cannot be written as a function of
\(\ell_v\), since that function would assign the same value to both states.
The difference of boundaries thus acquires a verifiable operator effect.

\subsection{Ambivalence weights and self-inertial response}
\label{vii:vii:sec:observador-ambivalencia}

For the general formula we consider positive readers
\(A_\partial,V_\partial\) and an energy reader \(Q\)
defined on the declared effective boundary domain.
A realization on
\(B_{\rm av}\subseteq B_{\rm ef}\) is made explicit below. Memory remains as a boundary
coordinate \(\mathsf M_\partial(\beta)\). The ambivalence
hinge provides the two orientations
\[
 r_{\rm av}=\frac{\pi_{\rm HMT}}{\varphi_{\rm HMT}^2},\qquad
 r_{\rm av}^{J}=\frac{\varphi_{\rm HMT}}{\pi_{\rm HMT}^2}.
\]
The superscript \(J\) expresses interchange of the generating pair.
\paragraph{Normalized realization of the boundary readers.}
The APP--TRIT--TPK modal incidence, developed in the regional
construction, provides the path envelope with matrix
\[
 F_{\rm av}=\begin{pmatrix}0&1\\1&1\end{pmatrix}.
\]
The positive vector \(v=(1,\varphi_{\rm HMT})^{\mathsf T}\)
satisfies \(F_{\rm av}v=\varphi_{\rm HMT}v\). Its normalized
transition and stationary weights are
\begin{equation}
 P_{ij}=\frac{(F_{\rm av})_{ij}v_j}{\varphi_{\rm HMT}v_i},
 \qquad
 P=\begin{pmatrix}0&1\\
 \varphi_{\rm HMT}^{-2}&\varphi_{\rm HMT}^{-1}\end{pmatrix},
 \qquad
 (\mu_{\rm ar},\mu_{\rm vol})
 =\frac{(1,\varphi_{\rm HMT}^{2})}{1+\varphi_{\rm HMT}^{2}}.
 \label{vii:vii:eq:lectores-parry}
\end{equation}
The rows of \(P\) sum to one by
\(\varphi^{-2}+\varphi^{-1}=1\); the stationary equation
\(\mu_{\rm ar}=\mu_{\rm vol}\varphi^{-2}\), together with
normalization, determines the stated weights. This Parry measure
weights the paths of the envelope; the chronological frequency
\((1/3,2/3)\) counts the instants of the modal calendar and retains
its distinct function.

The common memory-amplitude realization is defined on the sector
\(\mathcal X_{\rm av}\subset\mathcal X_\infty^{\rm enr}\) in which
the readers of the two sheets have the factorization
\begin{equation}
 \mathcal L_{\rm vol}(x)=\mu_{\rm vol}\mathcal M(x),\qquad
 \mathcal L_{\rm ar}(x)=\pi_{\rm HMT}\mu_{\rm ar}\mathcal M(x),
 \qquad \mathcal M(x)>0.
 \label{vii:vii:eq:lectores-amplitud-comun}
\end{equation}
Both values belong to the same positive memory-amplitude
line. The areal publication carries the regional mark
\(\pi_{\rm HMT}\) of the return. The dimensionless reading is obtained
by dividing both sections by \(\mathcal L_{\rm vol}(x)\).
Thus the sums of the weights are performed between quantities of the
same type; a possible realization as physical areas and volumes
is then composed with their respective dimensional transducers.

\begin{proposition}[Descent of the normalized readers and the energy clock]
\label{vii:vii:prop:lectores-frontera-explicitos}
On \(B_{\rm av}=b(\mathcal X_{\rm av})\), the rules
\begin{equation}
 A_\partial(b(x))
 =\frac{\mathcal L_{\rm ar}(x)}{\mathcal L_{\rm vol}(x)}
 =r_{\rm av},\qquad V_\partial(b(x))=1
 \label{vii:vii:eq:lectores-normalizados}
\end{equation}
define two positive readers independent of the representative.
Given an orientation \(a\) of the action section, the decision
energy of \eqref{vii:v:ant:eq:ii-kb-aut-energia} determines on states
\begin{equation}
 q_a(x)=\frac{h_a(x)}{108t_0(x)}
       =\frac{\hbar_a(x)}{t_P(x)}\in\mathcal L_E^+,
 \qquad t_P=\frac{54}{\pi_{\rm HMT}}t_0.
 \label{vii:vii:eq:lector-q-accion}
\end{equation}
We write
\[
 b_{\rm av}:=b|_{\mathcal X_{\rm av}}:
 \mathcal X_{\rm av}\longrightarrow B_{\rm av}.
\]
There exists a unique reader \(Q:B_{\rm av}\to\mathcal L_E^+\) such that
\(q_a|_{\mathcal X_{\rm av}}=Q\circ b_{\rm av}\)
if and only if, for any \(x,x'\in\mathcal X_{\rm av}\),
\begin{equation}
 b(x)=b(x')\ \Longrightarrow\
 \frac{h_a(x)}{t_0(x)}=\frac{h_a(x')}{t_0(x')}.
 \label{vii:vii:eq:descenso-reloj}
\end{equation}
In particular, a chart of fixed action and duration satisfies this
condition and publishes \(Q=h_a/(108t_0)\).
\end{proposition}
\begin{proof}
The quotient of \eqref{vii:vii:eq:lectores-amplitud-comun} cancels the
common amplitude and equals
\(\pi_{\rm HMT}\mu_{\rm ar}/\mu_{\rm vol}
=\pi_{\rm HMT}/\varphi_{\rm HMT}^2\).
The result is the same for all representatives of a fiber
of \(b_{\rm av}\) within \(\mathcal X_{\rm av}\) and is positive.
Dividing both readers by the same positive factor
preserves the weighted quotients in which they enter.
For the clock, the factorization criterion of
theorem~\ref{vii:vii:thm:observador-factorizacion}, applied to \(q_a\),
is exactly equivalent to \eqref{vii:vii:eq:descenso-reloj}. Fixed sections
satisfy it; so do variable sections whose
quotient descends through \(b_{\rm av}\). Uniqueness follows from the definition
of \(B_{\rm av}\) as the effective image.
\end{proof}

The value of the reader \(Q\) belongs to the energy line
\(\mathcal L_E=\mathcal L_{\rm act}\otimes\mathcal L_T^*\).
With \(E_1,E_2\) on that same line,
\(E_1E_2/Q^2\) is dimensionless and independent of the positive energy
basis. The symbol \(Q\) in this section exclusively designates
that clock energy; the quantity of heat in the Landauer
realization belongs to the corresponding thermal balance.
Common-amplitude factorization specifies a concrete
realization of the general readers. Every other domain preserves its
boundary maps and its own verification of descent.
Complete memory remains in the enriched history and in its
boundary publication, even though it cancels in the normalized quotient.


We define
\[
 w_A(\beta)=
 \frac{A_\partial(\beta)r_{\rm av}}
 {A_\partial(\beta)r_{\rm av}+V_\partial(\beta)r_{\rm av}^{J}},
 \qquad w_V(\beta)=1-w_A(\beta).
\]
In realization \eqref{vii:vii:eq:lectores-normalizados}, the weights
admit the exact evaluation
\begin{equation}
 w_A=\frac{r_{\rm av}^{\,2}}{r_{\rm av}^{\,2}+r_{\rm av}^{J}}
     =\frac{\pi_{\rm HMT}^{4}}
            {\pi_{\rm HMT}^{4}+\varphi_{\rm HMT}^{5}},\qquad
 w_V=\frac{\varphi_{\rm HMT}^{5}}
            {\pi_{\rm HMT}^{4}+\varphi_{\rm HMT}^{5}}.
 \label{vii:vii:eq:pesos-amplitud-comun}
\end{equation}
Substitution preserves both factors of the definition: the normalized
areal reader and its subsequent weighting by \(r_{\rm av}\).
Generation and regional marking of the return are performed once.
Multiplying numerator
and denominator by \(\varphi_{\rm HMT}^4\pi_{\rm HMT}^2\)
gives the last expression. These weights remain constant in the
common-amplitude sector. With fixed projectors and energies, a
variation of the response may come from the reader \(Q\);
the witness to global dependence additionally requires the states with
equal local restriction and different response of
\eqref{vii:vii:eq:observador-testigo-global}.


Let \(P_A,P_V\) be complementary orthogonal projections of a
response space \(\mathcal H_v\), and let \(E_1,E_2>0\) be
two internal energies fixed for comparison. The response is
\begin{equation}
 \overline I_v(\beta)=\frac{E_1E_2}{Q(\beta)^2}
       \bigl(w_A(\beta)P_A+w_V(\beta)P_V\bigr),
 \qquad I_v=\overline I_v\circ b.
 \label{vii:vii:eq:observador-respuesta-ponderada}
\end{equation}
Energies that are explicit readers of
\(\beta\) are also admitted; in that case they are preserved as such in the formula.

\begin{proposition}[Positivity and covariance of the response]
\label{vii:vii:prop:observador-respuesta-positiva}
The response \eqref{vii:vii:eq:observador-respuesta-ponderada} is
self-adjoint, positive definite and constant on the fibers of \(b\).
Its eigenvalues in the nonzero sectors are
\[
 \lambda_A=E_1E_2Q^{-2}w_A,\qquad
 \lambda_V=E_1E_2Q^{-2}w_V.
\]
The complete exchange
\((A_\partial,r_{\rm av},P_A)
 \leftrightarrow(V_\partial,r_{\rm av}^{J},P_V)\)
permutes the two summands and preserves the operator. A unitary
transport of the response space conjugates that operator and preserves
its spectrum with multiplicities.
\end{proposition}

\begin{proof}
Positivity of the readers gives \(0<w_A,w_V<1\). For
\(z=z_A+z_V\), with \(z_A=P_Az\), \(z_V=P_Vz\),
\[
 \langle z,\overline I_vz\rangle
 =\frac{E_1E_2}{Q^2}
   (w_A\|z_A\|^2+w_V\|z_V\|^2)>0
\]
if \(z\neq0\). The orthogonal decomposition proves the spectral
expressions. Exclusive dependence on \(\beta\) gives constancy
on the fibers and factorization. Exchanging the triples
exchanges numerators, weights and projectors at once, so it
preserves the sum. For a unitary transport \(U\), substitution
\(P_A\mapsto UP_AU^*\), \(P_V\mapsto UP_VU^*\) produces
\(U\overline I_vU^*\), with the same spectrum.
\end{proof}

An exchange implemented within the same space by a unitary
sending \(P_A\) to \(P_V\) requires the two sectors to have
equal dimension. Exchanging the two triples as labels of the
formula is valid with each sector's own dimension.

Condition \eqref{vii:vii:eq:observador-testigo-global} requires comparing
the responses. For example, simultaneously multiplying
\(A_\partial,V_\partial\) by the same positive factor preserves
both weights; if \(E_1,E_2,Q\) are also preserved, the response
remains identical. In contrast, with equal weights and energies, replacing
\(Q\) by \(qQ\), \(q>0\), multiplies the response by \(q^{-2}\).
If two states with equal local restriction realize that change with
\(q\neq1\), they provide the global witness. These are operator
separation conditions; existence of the states is checked
in the boundary domain being realized.

\subsection{Connection of the joint state and acceleration reading}
\label{vii:vii:sec:observador-autoinercia}

In a differentiable realization of the joint state
\(\Gamma=(x,\mathcal O,m)\), let \(N\) be its realization space,
\(\nabla^{\rm tot}\) its connection and \(\pi_S:N\to N_S\) a
differentiable projection onto the subsystem, equipped with connection
\(\nabla^S\). The connection and projection preserve the provenance
of the transports and readers they realize. The observed frame is
obtained through
\[
 \mathfrak F_{\mathcal O}(\Gamma)
 =\operatorname{Res}_{\mathcal O}(\Gamma,\nabla^{\rm tot}).
\]
For a twice-differentiable trajectory \(\Gamma(t)\),
write \(\gamma_S=\pi_S\circ\Gamma\).

\begin{definition}[Self-inertia and projection defect]
\label{vii:vii:def:observador-autoinercia}
A trajectory is self-inertial with respect to the declared connection
when \(\nabla^{\rm tot}_{\dot\Gamma}\dot\Gamma=0\).
The defect of its projection is
\begin{equation}
 \mathcal K_{\pi_S}(\dot\Gamma,\dot\Gamma)
 =\nabla^S_{\dot\gamma_S}\dot\gamma_S
  -d\pi_S\bigl(\nabla^{\rm tot}_{\dot\Gamma}\dot\Gamma\bigr).
 \label{vii:vii:eq:observador-defecto-proyeccion}
\end{equation}
\end{definition}

\begin{proposition}[Subsystem acceleration]
\label{vii:vii:prop:observador-aceleracion}
For a self-inertial trajectory,
\[
 a_S=\nabla^S_{\dot\gamma_S}\dot\gamma_S
     =\mathcal K_{\pi_S}(\dot\Gamma,\dot\Gamma).
\]
In coordinates \(x^B\) on \(N\) and \(y^a\) on \(N_S\), the
defect is quadratic in velocity, with components
\begin{equation}
 (\mathcal K_{\pi_S})^a{}_{BC}
 =\partial_B\partial_C\pi_S^a
  +(\Gamma^S)^a{}_{bc}\,
        \partial_B\pi_S^b\partial_C\pi_S^c
  -\partial_D\pi_S^a(\Gamma^{\rm tot})^D{}_{BC}.
 \label{vii:vii:eq:observador-defecto-coordenadas}
\end{equation}
\end{proposition}

\begin{proof}
The chain rule gives
\(\ddot\gamma_S^a=\partial_B\pi_S^a\ddot\Gamma^B
 +\partial_B\partial_C\pi_S^a\dot\Gamma^B\dot\Gamma^C\).
Adding the contribution of \(\nabla^S\) and subtracting the image of
the total acceleration, the terms containing \(\ddot\Gamma\)
cancel. The remaining ones are
\eqref{vii:vii:eq:observador-defecto-coordenadas} contracted with
\(\dot\Gamma^B\dot\Gamma^C\). When total acceleration vanishes,
the stated equality results.
\end{proof}

The geometric response factors through \(b\) when the connection,
projection and kinematic datum used in evaluation are
determined by \(\beta=b(x)\), or when those kinematic data are
kept fixed in the comparison. Under those conditions,
\eqref{vii:vii:eq:observador-defecto-coordenadas} gives the same result
for any representatives of a fiber;
theorem~\ref{vii:vii:thm:observador-factorizacion} constructs its unique
factor on \(B_{\rm ef}\). If the kinematic datum varies, it
forms part of the domain of the response being compared.

The self-inertial formulation describes a geodesic trajectory of the
joint state and an effective acceleration of its subsystem due
to the projection defect. The gravitational realization then identifies
that connection, the frame and the response with their physical
observables. Observer covariance, memory inscription and
factorization through the boundary thus retain complementary functions:
the first determines the reading class, the second retains the
genealogy and the third specifies the relational dependence of the
response.
